\section{Introduction}\label{sec:introduction}

% Mathematical job:
%   - State the integrated thesis of Foundations II (stratified closure
%     calculus, not flat six-generator algebra).
%   - Introduce the three co-equal pillars: primitive role architecture,
%     admissibility / meta-theory, and the scoped exact-six program.
%   - State the terminal theorem at scoped strength and its hypotheses.
%   - State the paper-wide nonclaims up front.
%
% Governing sources:
%   vision.md
%   workspace/meta/lens_lab/papers/paper3/paper3_outline.md
%   workspace/meta/lens_lab/papers/paper3/theorem_claims.md
%   workspace/meta/lens_lab/papers/paper3/scope_and_nonclaims.md
%   paper/orientation_memo.md (sections 1, 2, 6)
%
% Overclaim risks:
%   - Writing as a Paper-3 report; centering the exact-six terminus.
%   - Importing programmatic language from vision.md.
%   - Dropping the scope qualifier from the main theorem.

Six Birds Foundations I~\citep{Tsiokos2026SixBirdsFoundations}
introduced a stratified closure-calculus
perspective~\citep{Anderson1972} on finite theory packages,
\(\mathcal{T} = (Z, f, \Sigma_f, E, \mathcal{A})\), in which theories
arise as fixed points of idempotent
operators~\citep{DaveyPriestley2002} and six primitive closure roles
\Pone{}, \ldots, \Psix{} emerge from composability combined with
limited-interface observation. The present paper extends
that calculus into a theorem-level framework. It fixes the primitive
role architecture, specifies the admissibility and meta-theory regime
under which role talk is licensed, and proves a scoped exact-six
terminal result over a precisely delimited domain. The exact-six
theorem is the consequence of that stack, not its whole content.

\subsection{The three parts}\label{subsec:three-pillars}

Each of the three parts below supplies structure the other two do not.
The \emph{primitive role architecture}
(Section~\ref{sec:primitive-role-architecture}) fixes the six role
meanings of \Pone{}--\Psix{} at their corrected alignments, records
the common behavioral, verification, and structural-categorical level
pattern, enumerates the local boundary distinctions under scoped
bookkeeping, and states the selected-lift atlas vocabulary. The
\emph{admissibility and meta-theory regime}
(Section~\ref{sec:admissibility-meta-theory}) elevates honest
bookkeeping from a per-description condition to a meta-theoretic
discipline, replaces flat primitive independence by typed
non-collapse, and records level, forgetting, threshold, visibility,
and empirical-bridge theorem schemas. The \emph{exact-six theorem
chain}
(Sections~\ref{sec:fatcd-domain}--\ref{sec:integrated-stress})
quantifies the exact-six question precisely over the domain \FATCD{}
of finite audited typed closure descriptions and proves the terminal
result via lower-bound witnesses, an upper-bound classification with
corrected role-projection alignments, a decomposition/exhaustion
theorem, and an integrated stress certification.

\subsection{The terminal theorem, informally}\label{subsec:terminal-informal}

A finite audited typed closure description \(D\) is a finitely
presented collection of typed raw records and annotations subject to
finiteness, closure, and honest-bookkeeping conditions. The role
vocabulary \Pone{}--\Psix{} enters \(D\) not as a constraint on
admissibility but as a family of derived projections \(\pi_i(D)\),
each available when threshold, witness, level, and audit data are
attached to raw features. The terminal theorem
(Theorem~\ref{thm:scoped-exact-six}) states that for every
\(D \in \FATCD{}\) there exists a well-formed
decomposition/exhaustion record \(\Dec{D}\) carrying exactly the six
typed closure-role channels \Pone{}, \ldots, \Psix{}, each in exactly
one of five statuses, with all residual features covered by an
explicit classification and no seventh irreducible role required in
the covered scope. The claim is about role channels, not active
projections, and the decomposition is existential, not unique.

\subsection{What the paper does and does not claim}\label{subsec:intro-nonclaims}

The terminal theorem is scoped exact-six, not unrestricted exact-six.
It makes none of the following stronger claims:
\begin{itemize}[topsep=2pt,itemsep=1pt]
  \item unrestricted exact-six, or any claim about all emergence
    phenomena;
  \item that every \(D \in \FATCD{}\) activates all six roles or
    carries six active nontrivial witnesses;
  \item that the decomposition/exhaustion record \(\Dec{D}\) is unique
    or canonical;
  \item a full six-primitives algebra or global primitive
    independence;
  \item empirical realization;
  \item internal instrument-family completeness;
  \item any statement about broader domains without further work.
\end{itemize}
These are scope boundaries, collected in
Section~\ref{sec:limitations-nonclaims} and enforced locally by the
theorem-chain results that would otherwise overread them.

\subsection{Organization}\label{subsec:organization}

After this introduction, Section~\ref{sec:primitive-role-architecture}
fixes the primitive role meanings, level trichotomy, boundary
distinctions, and selected-lift atlas. Section~\ref{sec:admissibility-meta-theory}
states the seven-schema admissibility regime. Section~\ref{sec:fatcd-domain}
defines \FATCD{} and formulates the exact-six question over it.
Section~\ref{sec:lower-bounds} proves the six-role lower bounds.
Section~\ref{sec:upper-bound-classification} proves the repaired
upper-bound classification theorem, closing the probability,
causality, and agency threats within the covered scope.
Section~\ref{sec:decomposition-exhaustion} proves decomposition
existence, six-channel exactness, and residual exhaustion, and records
the channel-versus-activation discipline and non-uniqueness.
Section~\ref{sec:integrated-stress} records the anti-overclaim stress
families and the integrated stress certification.
Section~\ref{sec:scoped-exact-six} states and proves the terminal
theorem. Sections~\ref{sec:limitations-nonclaims}
and~\ref{sec:future-work} record scope boundaries and future
directions. Appendix~\ref{app:api-provenance} records the API
provenance of the theorem-chain artifacts.
